%Chargement des paquets
\usepackage{amsmath}
\usepackage{amsfonts}
\usepackage{amssymb}
\usepackage{enumerate}
\usepackage{mathtools}
\usepackage{bbm}
\usepackage{xparse, etoolbox}
\usepackage{enumerate}
\usepackage{enumitem}
\usepackage{mathabx}
\usepackage{babel}[french]

%environment
\usepackage{amsthm}
\usepackage{keytheorems}
\usepackage{hyperref} 

%Interface théorème
\newkeytheorem{lem}[
	name = Lemme
]

\newkeytheorem{prop}[
	name = Proposition
]

\newkeytheorem{thm}[
	name = Théorème
]

\newkeytheorem{coro}[
	name = Corollaire
]

\renewcommand*{\proofname}{Démonstration}

\theoremstyle{definition}
\newtheorem{defn}{Définition}

\theoremstyle{definition}
\newtheorem*{nota}{Notation}

\theoremstyle{definition}
\newtheorem*{limi}{Liminaire}

\theoremstyle{remark}
\newtheorem{rmq}{Remarque}

\theoremstyle{plain}
\newtheorem*{fait}{Fait}


%Ensembles de nombres
\newcommand{\N} {\mathbb{N}}
\newcommand{\Ne}{\N^\ast}
\newcommand{\Z} {\mathbb{Z}}
\newcommand{\Q} {\mathbb{Q}}
\newcommand{\R} {\mathbb{R}}
\newcommand{\Rb}{\overline{\mathbb{R}}}
\newcommand{\Rp}{\R_+}
\newcommand{\Rm}{\R_-}
\newcommand{\K} {\mathbb{K}}
\newcommand{\Ket} {\mathbb{K^{\ast}}}
\newcommand{\Cx}{\mathbb{C}}

%Ensembles, tribus
\newcommand{\A}{\mathbb{A}}
\renewcommand{\P}{\mathcal{P}}
\newcommand{\T}{\mathcal{T}}
\newcommand{\F}{\mathcal{F}}
\newcommand{\C} {\mathcal{C}}
\newcommand{\ouv}{\mathcal{O}}
\newcommand{\B}{\mathcal{B}}
\newcommand{\M}{\mathcal{M}}
\newcommand{\etag}{\mathcal{E}}
\newcommand{\etagOp}{\etag^{+}(\Omega)}
\newcommand{\etagO}{\etag(\Omega)}
%%Lp avant quotientage
\newcommand{\Lic}{\mathcal{L}^1}
\newcommand{\Lpc}{\mathcal{L}^p}
\newcommand{\Linfc}{\mathcal{L}^{\infty}}
\newcommand{\Lifull}{\Lic(\Omega, \B, m)}
\newcommand{\Lpfull}{\Lpc(\Omega, \B, m)}
\newcommand{\Linffull}{\Linfc(\Omega, \B, m)}
%%Lp après quotientage
\newcommand{\Li}{L^1}
\newcommand{\Ld}{L^2}
\newcommand{\Lp}{L^p}
\newcommand{\Linf}{L^{\infty}}
\newcommand{\Listd}{\Li(\Omega, m)} 
\newcommand{\Ldstd}{\Ld(\Omega, m)} 
\newcommand{\Lpstd}{\Lp(\Omega, m)} 

%Quelques opérateurs
\DeclareMathOperator{\codim}{codim}
\DeclareMathOperator{\rg}{rg}
\DeclareMathOperator{\tr}{tr}
\DeclareMathOperator{\vect}{Vect}
\DeclareMathOperator{\ima}{Im}
\DeclareMathOperator{\id}{Id}
\DeclareMathOperator{\sign}{sign}
\DeclareMathOperator{\ext}{ext}
\newcommand{\ind}{\mathbbm{1}}
\newcommand*{\spr}[2]{\left(\,#1\,\middle|\,#2\,\right)}
\newcommand{\interior}[1]{%
  {\kern0pt#1}^{\mathrm{o}}%
}
\DeclareMathOperator{\card}{card}
\DeclareMathOperator{\ppcm}{ppcm}

%Commandes de pure flemme
\newcommand{\veps}{\varepsilon}
\newcommand{\intab}{\int_{a}^{b}}
\newcommand{\intR}{\int_{-\infty}^{\infty}}
\newcommand{\intinf}{\int_{- \infty}^{+ \infty}}
\newcommand{\equi}{\Leftrightarrow}
\newcommand{\Kev}{$\K$-espace vectoriel }
\newcommand{\Kevs}{$\K$-espaces vectoriels }
\newcommand{\Rev}{$\R$-espace vectoriel }
\newcommand{\Revs}{$\R$-espaces vectoriels }
\newcommand{\Cev}{$\Cx$-espace vectoriel }
\newcommand{\Cevs}{$\Cx$-espaces vectoriels }
\newcommand{\evn}{(E, \norm{.})}
\newcommand{\interf}[2]{[#1,#2]}
\newcommand{\limb}{\lim\limits_}
\newcommand{\lebn}{\lambda_n}
\newcommand{\pp}{\text{~presque partout}}
\newcommand{\derivp}[2]{\frac{\partial #1}{\partial #2}}
\newcommand{\famiu}{(u_i)_{i \in I}}
\newcommand{\sev}{\text{~s.e.v.}}
\newcommand{\Mnp}{M_{n,p}(\K)}
\newcommand{\Mn}{M_{n}(\K)}
\newcommand{\tq}{\text{~t.q.~}}
\newcommand{\mq}{~montrer~que~}
\newcommand{\et}{\quad \text{et} \quad}
\newcommand{\ou}{\text{~ou~}}
\newcommand{\Kx}{\K[X]}


%Commandes de gars chelou
\renewcommand{\subset}{\subseteq}

%Commande restriction, ty duckduckgo
\def\restr#1#2{\mathchoice
              {\setbox1\hbox{${\displaystyle #1}_{\scriptstyle #2}$}
              \restraux{#1}{#2}}
              {\setbox1\hbox{${\textstyle #1}_{\scriptstyle #2}$}
              \restraux{#1}{#2}}
              {\setbox1\hbox{${\scriptstyle #1}_{\scriptscriptstyle #2}$}
              \restraux{#1}{#2}}
              {\setbox1\hbox{${\scriptscriptstyle #1}_{\scriptscriptstyle #2}$}
              \restraux{#1}{#2}}}
\def\restraux#1#2{{#1\,\smash{\vrule height .8\ht1 depth .85\dp1}}_{\,#2}} 

%Quelques normes
\DeclarePairedDelimiter\abs{\lvert}{\rvert}
\DeclarePairedDelimiter\labs{\bigl\lvert}{\bigr\rvert}
\DeclarePairedDelimiter\norm{\lVert}{\rVert}
\DeclarePairedDelimiterXPP\normi[1]{}\lVert\rVert{_1}{\ifblank{#1}{\:\cdot\:}{#1}}
\DeclarePairedDelimiterXPP\normd[1]{}\lVert\rVert{_2}{\ifblank{#1}{\:\cdot\:}{#1}}
\DeclarePairedDelimiterXPP\normp[1]{}\lVert\rVert{_p}{\ifblank{#1}{\:\cdot\:}{#1}}
\DeclarePairedDelimiterXPP\normq[1]{}\lVert\rVert{_q}{\ifblank{#1}{\:\cdot\:}{#1}}
\DeclarePairedDelimiterXPP\norminf[1]{}\lVert\rVert{_\infty}{\ifblank{#1}{\:\cdot\:}{#1}}

%Bracket en angle
\DeclarePairedDelimiter\angl{\langle}{\rangle}

%Set, parenthèses
\newcommand{\set}[1]{\{ #1 \}}
\newcommand{\bigset}[1]{\bigl\{ #1 \bigr\}}
\newcommand{\Bigset}[1]{\Bigl\{ #1 \Bigr\}}
\newcommand{\bigpar}[1]{\bigl( #1 \bigr)}
\newcommand{\Bigpar}[1]{\Bigl( #1 \Bigr)}

%Opitmisation des opérateurs de comparaison
\DeclarePairedDelimiter\tio{o (}{)}
\DeclarePairedDelimiter\gro{O (}{)}

\newcommand{\Ee}{E^{\ast}}
\newcommand{\zero}{\set{0}}

\pagestyle{fancy}

\fancyhead[C]{}
\fancyhead[R]{}

\begin{document}

\underline{Source :} 
\begin{itemize}
\item Grifone - Algèbre linéaire - page 289 ;
\item Rombaldi - Algèbre - page 476 (back-up).
\end{itemize}

\underline{Leçons :}
\begin{itemize}
\item  159 : Formes linéaires et dualité en dimension finie. Exemples et applications.
\item  170 : Formes quadratiques sur un espace vectoriel de dimension finie. Orthogonalité. Applications.
\item  171 : Formes quadratiques réelles. Coniques. Exemples et applications.
\item  148 : Dimension d’un espace vectoriel (on se limitera au cas de la dimension finie). Rang. Exemples et applications.
\end{itemize}

\begin{thm}
Soit $q$ une forme quadratique sur $E$. Il existe une base $q$-orthogonale de $E$.
\end{thm}

\begin{proof}
On démontre ce théorème par récurrence sur la dimension $n$ de $E$.
Le théorème est vrai pour $n=1$. Supposons le vrai à l'ordre $n-1 \geq 1$ et démontrons le pour $E$ de dimension $n$.
Le résultat est trivial pour $q$ identiquement nulle, supposons donc que $q$ ne soit pas telle et prenons $v \in E$ tel que $q(v) \neq 0$.
On pose :
\[ F = \set{x \in E ; s(x,v) = 0 } , \]
où $s$ est la forme polaire de $q$. 
On pose $\varphi \in \Ee$ la forme linéaire définie par $\forall x \in E, \varphi(x) = s(x,v)$.
Alors $F = \ker(\varphi)$ et comme $\varphi(v) = s(v,v) = q(v) \neq 0$, on déduit que $F$ est un hyperplan de $E$.
On a donc $E = F \oplus D$ et comme $v \not \in F$, on a $E = F \oplus \vect(v)$.
On peut appliquer l'hypothèse de récurrence sur $F$, obtenir une base $q$-orthogonale de $F$ et la compléter avec le vecteur $v$
pour démontrer que le résultat est héréditaire, ce qui clos la récurrence.
\end{proof}

\begin{rmq}
Dans une base orthogonale, la matrice $M$ de $q$ est de la forme :
\[ M = \begin{pmatrix}
a_1 & & & & & 0 \\
& \ddots & & & & \\
& & a_r & & & \\
& & & 0 & & \\
& & & & \ddots & \\
& & & & & 0
\end{pmatrix} \]
\end{rmq}

\begin{thm}[de Sylvester]
On suppose que $\K = \R$.
Soit $q$ une forme quadratique sur $E$, il existe alors une base $(e_1, \dots, e_n)$ de $E$ telle que :
\[ \forall x \in E, x = \sum x_i e_i, \quad q(x) = x_1^2 + \dots + x_p^2 - x_{p+1}^2 - \dots x_r^2 , \]
c'est-à-dire que la matrice $M$ de $q$ dans cette base est de la forme :
\[ M = \begin{pmatrix}
\begin{matrix} 1&&\\ &\ddots& \\&&1 \end{matrix} &&0 \\
& \begin{matrix} -1&&\\ &\ddots& \\&&-1 \end{matrix} & \\
0 && \begin{matrix} 0&&\\ &\ddots& \\&&0 \end{matrix}
\end{pmatrix} \]
où $r = \rg(q)$ et $p$ est un entier ne dépendant que de $q$ (et non de la base choisie).
Le couple $(p,r-p)$ est noté $\sign(q)$ et est appelé signature de $q$.
\end{thm}

\begin{proof}
Soit $(e_1, \dots, e_n)$ une base $q$-orthogonale de $E$. Si $x = \sum y_i e_i$, on a alors :
\[ q(x) = a_1 y_1^2 + \dots a_r y_r^2 , \]
avec les $a_i$ des réels et $r \geq n$ un entier naturel.
Quitte a changer la numérotation, on peut supposer que $a_1, \dots, a_p >0$ et $a_{p+1}, \dots, a_r < 0$, et on écrit
\begin{align*}
q(x) & = (\sqrt{a_1} y_1)^2 + \dots + (\sqrt{a_p} y_p)^2 - (\sqrt{-a_{p+1}} y_{p+1})^2 + \dots + (\sqrt{-a_r} y_r)^2 \\
& = x_1^2 + \dots + x_p^2 - x_{p+1}^2 - \dots - x_r^2 , 
\end{align*}
avec $x_i = \sqrt{a_i} y_i, i \in \set{1, \dots, p}$ et $x_j = -\sqrt{a_j} y_j, j \in \set{p+1, \dots, r}$,
(et la matrice de $q$ a la forme cherchée dans la base 
$\left ( \dfrac1{\sqrt{a_1}} e_1, \dots, \dfrac1{\sqrt{a_p}} e_p, \dfrac1{\sqrt{a_{p+1}}} e_{p+1}, \dots, \dfrac1{\sqrt{a_{r}}} e_{r}, e_r, \dots, e_n
\right )$ ).
Il reste à montrer que l'entier $p$ ne dépend pas du choix de la base.
Considérons donc une seconde base $q$-orthogonale notée $(f_1, \dots, f_n)$.
Dans cette base, on a $x=\sum_i y'_i f_i$ et en appliquant le même raisonnement que précédemment on peut écrire :
\[ q(x) =  {x'_1}^2 + \dots + {x'_{p'}}^2 - {x'_{p'+1}}^2 - \dots - {x'_r}^2 . \]
On va maintenant considérer les quatre \sev de $E$ suivant :
\[ \begin{matrix}
F = \vect(e_1, \dots, e_p) & \qquad & F' = \vect(e'1, \dots, e'_{p'}) \\
G = \vect(e_{p+1}, \dots, e_n) & \qquad & G' = \vect(e'_{p'+1}, \dots, e'_n) . 
\end{matrix} \]
Considérons $x \in F \cap G'$, alors $q(x) \geq 0$ car $x \in F$ et $q(x) \leq 0$ car $x \in G$.
Donc $F \cap G' = \zero$, ce qui signifie que $F et G'$ sont en somme directe dans $E$.
Nécessairement, on a $p+n-p' \leq n$ et donc $p \leq p'$.
En prenant maintenant $x' \in F' \cap G$, on conduit le même raisonnement et on trouve $p' \leq p$.
Finalement $p=p'$ et on a démontré le résultat.
\end{proof}

\begin{coro}
On suppose toujours que $\K = \R$. Soit $q$ une forme quadratique sur $E$, on a alors le schéma d'équivalences suivant :
\[ \begin{matrix}
q \text{ est définie positive } & \iff & \sign(q) = (n,0) \\
\big\Updownarrow & & \big\Updownarrow \\
E \text{ euclidien } & & \text{ il existe des bases orthonormées} \\
q \text{ est définie négative } & \iff & \sign(q) = (0,n) \\
q \text{ est non dégénérée } & \iff & \sign(q) = (p,n-p) \\
\end{matrix} \]
\end{coro}


\end{document}